\documentclass[10pt,letterpaper]{article}
\usepackage[T1]{fontenc}
\usepackage{lmodern,amsmath,amssymb,amsthm,microtype}
\usepackage[margin=0.84in]{geometry}
\usepackage[hidelinks]{hyperref}
\providecommand{\tightlist}{\setlength{\itemsep}{0pt}\setlength{\parskip}{0pt}}
\DeclareUnicodeCharacter{2264}{\ensuremath{\le}}
\DeclareUnicodeCharacter{2265}{\ensuremath{\ge}}
\DeclareUnicodeCharacter{2260}{\ensuremath{\ne}}
\DeclareUnicodeCharacter{2212}{\ensuremath{-}}

\DeclareUnicodeCharacter{2211}{\ensuremath{\sum}}

\title{Six grid additions to \(G_{449}\): a universal hub with one base neighbor}
\author{Aafi Khan Malek\\ \small Independent researcher, Ahmedabad, India\\
\small\href{mailto:aafimalek2032@gmail.com}{aafimalek2032@gmail.com}}
\date{28 November 2026}
\begin{document}
\maketitle
\textbf{Theorem.} Add six distinct points in the denominator-seven
cyclotomic grid to the Heule--Przybocki graph \(G_{449}\). If exactly
one added point has one original neighbor and it is adjacent to all
five other added points, then the induced extension is five-colorable.
The proof uses the previously verified
\href{https://github.com/Aafimalek/novel-discoveries/tree/234e0ce194fdfa544d4b551ccea9898b2239dc69/hadwiger-nelson/verification/g449/prime-seven-catalogue/five-point-extensions}{five-point extension theorem}.
This is not a new chromatic bound for the plane or a Lean theorem.

\section{Derivation}\label{derivation}

Assume the prior theorem that any five points can be added to G449 with
five colors. For a new vertex \(v\) let \(b(v)\) be its number of
original neighbors and \(d_H(v)\) its degree among the six additions. A
counterexample must have \[
\delta_{\rm add}=\min_v(b(v)+d_H(v))\ge5
\] by deleting and greedily restoring a vertex of smaller total degree.

In the checked subcase the unique vertex \(h\) with \(b(h)=1\) is
adjacent to all five other new vertices. All other vertices have
\(b\ge2\): another degree-one vertex is outside the case, and a
degree-zero vertex would also be universal in a putative counterexample
and share at least three unit neighbors with \(h\), impossible. The
7,056 possible rim points have \(b\ge2\); enumerate the 546 exact signed
grid steps from each, in total 3,852,576 incidences. Retain only hub
points with exactly one original neighbor, storing its original index
from the complete coordinate lookup. All incidences at a hub enter the
same shard. The earlier independent degree-one incidence histogram
predicts 109,683 retained incidences on 30,233 hubs, 9,394 hubs with at
least five rims, and 1,350,118 raw five-rim choices. The collector
checks these aggregate counts.

On the hub's unit circle the available rim graph has maximum degree two
and components of at most six vertices (unit chords move by
\(\pm60^\circ\)). A selected rim requires \(b(v)+d_{\rm rim}(v)\ge4\).
Peel any rim vertex that fails even in the full remaining rim graph;
then enumerate eligible subsets in each component and combine them to
size five, as in the already checked degree-zero universal-hub verifier.
For each proper base palette, forbid the hub's unique original
neighbor's color. For each other hub color in turn, forbid that color at
all rim vertices in addition to the colors of their respective base
neighbors. Run exact finite list search. No solver is called, and a bank
residual is only an uncolored configuration under those palettes, not a
counterexample.

\section{Exact verification}
All 32 disjoint shards completed, covering 1,350,118 raw five-rim
choices. There were 5,208 configurations after the minimum-degree
test, all five-colored using the existing palette bank. The first
18 palettes suffice. A fresh independent C++20 replay matched the
original frozen inputs and every shard receipt exactly, with no SAT
solve. The nonuniversal hub branch is separate.

\end{document}
